\section{Spacecraft Architecture and Ground Segment}\label{spacecraft-architecture}
The ANTARA spacecraft is a 3U CubeSat with the on-board computer as the main command and operations unit. The subsystem design currently uses preliminary COTS components along with the in-house scientific payload. Detailed specifications, operating estimates, and the original component tables are provided in Appendix~\ref{app:source-03}.

\subsection{Electrical Power System (EPS)}\label{electrical-power-system-eps}
The spacecraft uses one fixed solar panel on the +Z face and two panels stowed on the -Y face that deploy toward the +Z direction. The mission power studies use a 25~Wh battery planning capacity. The proposed ISISPACE ICEPS2 Type A has a manufacturer-rated 28.8~Wh battery-pack configuration; the mission simulation is assessed at 25~Wh. The original RAAN trade study reports average generated power of 11.05~W and 11.40~W for $0^\circ$ and $30^\circ$ RAAN respectively. These are study-case predictions, not a guarantee at any launch RAAN. The power studies were performed using end-of-life assumptions; their margins require acceptance against an agreed common run baseline.

The EOL power simulations provide a provisional sizing basis. Solar active area, conversion losses and worst-case EOL margin must be explicitly tied to an approved simulation configuration before these results are reported as closed budgets (Appendix~\ref{app:detailed-budgets}).

\subsection{Communications System (COMMS)}\label{communications-system-comms}
The preliminary communications system uses the GomSpace NanoCom AX100 and directional and omnidirectional spacecraft antennas. The planned primary ground station is at Goa; a crossed-Yagi ground antenna is being considered. Ground-station visibility is assessed using $10^\circ$ nominal and $15^\circ$ conservative elevation masks. A secondary station is an option if the primary station cannot support the required operations.

The final downlink rate and storage requirement cannot yet be fixed, because the payload packet format and science-data generation are not finalized. The original 15--18~kbps and simulator 19.2~kbps values are study assumptions, not a demonstrated data-return capacity. The ground-station contact case and data-rate calculations are described in Appendix~\ref{app:orbit-rf}.

\subsection{Scientific Payload}\label{payload-subsystem}
The current FPDU concept uses five $6\times6\times14$~mm main scintillator blocks in series along X, plus one $6\times6\times2$~mm front scintillator. Particles enter the $+X$ face and travel inward toward $-X$. Main-block SiPMs attach to $\pm Y$ (ten sensors), with one further front SiPM: eleven total. Teflon coats $\pm Z$ at 0.11~mm per face, giving 6.22~mm coated Z dimension. EJ-550 grease provides optical coupling. These are \textbf{CONFIRMED DESIGN INPUT}; material, inter-block optical isolation and packaged CAD remain open. The historical silicon entrance is removed; a rear scintillator is optional and unselected. The proton-energy requirement remains 5--100~MeV.

The proposed chain is scintillator, SiPM, ASIC, ADC, MCU or FPGA, classification, reconstruction, 64-bin histogram, and storage/downlink. ASIC shaping is 20~ns. Coincidence-window references are 200~ns nominal, at least 250~ns including allowance, and approximately 300~ns preferred. A separate stated 20~ns post-shaping interval has unresolved meaning. Qualified dead time, ADC latency and processing latency remain unknown. MCU is preferred only if the entire architecture meets timing, throughput, buffering and power needs; no processor is selected.

The historical \textbf{PRELIMINARY ESTIMATE} of 400~Hz assumed protons only, 250~$\mu$m (0.25~mm) aluminium shielding, geometric factor 1~cm$^2$~sr and 100\% detection efficiency; its generating calculation remains unidentified. Supplied per-cell rates supersede that rounded estimate. Reproduced proton maxima are 358.920~Hz (mean) and 975.659~Hz (99th-percentile map), using isotropic conversion, $G=1$~cm$^2$~sr, unit efficiency and an integral threshold of 5.5~MeV. These 2026-01-01, 600~km maps are approved preliminary references for 2029 planning. Their rate script does not model shielding or accepted detector response. Electron rates also contribute to electronics design studies. Appendix~\ref{app:payload-detail} gives the full assumptions, diagrams, timing trade and data contract; actual maximum trigger throughput is unestablished.

\subsection{Attitude Determination and Control System (ADCS)}\label{attitude-determination-and-control-system-adcs}
ANTARA intends to use an integrated reaction-wheel and magnetorquer attitude-control system; the earlier CMG description is superseded. The supplied TensorTech ADCS-10m databook describes a CMG-equipped unit and therefore does not confirm the intended reaction-wheel plus magnetorquer architecture. The ADCS vendor/product selection remains open. The main attitude functions are initial detumbling, Sun Pointing, payload +X pointing during SAA operations, +Y relay pointing toward the ground station, and Safe Mode recovery. Manufacturer limits and exact actuator configuration will be entered after the supplier data are reconciled with the selected model. The preliminary pointing goal remains subject to simulated and hardware verification. Mode transitions are given in Appendix~\ref{app:operational-logic}.

\subsection{On-Board Computer (OBC)}\label{on-board-computer-obc}
The source design uses the GomSpace NanoMind A3200 as the preliminary on-board computer. The OBC is responsible for mission sequencing, spacecraft state monitoring, command execution, payload control, data handling, and communications scheduling. The source document gives approximately 0.17~W as the average OBC power consumption for its stated operating conditions.

\subsection{Spacecraft Coordinate System}\label{spacecraft-coordinate-system}
The spacecraft body frame uses +X as the RAM face and payload boresight, +Y as the directional antenna face, +Z as the fixed solar-panel face, -Y as the deployable-panel packed face, -Z as the omnidirectional communications antenna face, and -X as the deployment face.

\subsection{Thermal and Structural Configuration}\label{sec:thermal-architecture}
Thermal and structural analyses are part of the MCR work, but the provided document does not yet establish qualified spacecraft temperature limits, verified mechanical modes, or completed environmental qualification. The supporting model sections are reserved in Appendix~\ref{app:thermal-structural}; their conclusions will be included here when validated results are available.

\subsection{Ground Segment and Mission Operations}\label{ground-segment-mission-operations}
The primary ground station is planned at the Goa campus, approximately $15^\circ$N, $74^\circ$E. A second station is a contingency only if the primary station cannot meet operational requirements. The ground segment provides tracking, telemetry reception, commanding, and science-data retrieval. The operational sequence uses Cold and Dark, Detumble, Sun Pointing, Detection, Relay, and Safe Mode, with Slewing as an attitude-transition state in the simulation. The full state-machine logic is given in Appendix~\ref{app:operational-logic}.

\subsection{Outstanding MCR Evidence}
\paragraph{Design maturity.} The MCR retains the reaction-wheel/magnetorquer architecture, selected solar-panel geometry and preliminary RF concept as a feasible spacecraft configuration. Detector response, finalized ADCS supplier interface, RF throughput and structural/thermal component compliance remain explicit PDR verification tasks rather than unspecified blanks.
